\documentclass[11pt,letterpaper]{article}
\usepackage{rf}

\addbibresource{references.bib}

\begin{document}

\begin{center}
    {\LARGE Symmetric Fixed Points of the Ricci Flow} \\[6pt]
    {\large \textbf{GPT Proof Write-Up}}
    \[
    \begin{array}{ll}
    \text{Author:}  &\text{Leo Lee} \\
    \text{Date:} &\text{\today} \\
    \end{array}
    \]
\end{center}

\vspace{-6pt}

\section{Introduction}
To find $(\SO(3) \times \SO(2))$-invariant Ricci solitons given by $(M, g)$ and a function $f: (\R^2 \times S^2) \to \R$, we will study the following ODE system \cite[Eqs.~(2.1)--(2.3)]{Rajan}
\begin{align}
    f'' &= \frac{a''}{a} + 2 \frac{b''}{b} - 1, \label{eq:rajan21} \\
    a'' &= -2 \frac{a'b'}{b} + a'f' + a, \label{eq:rajan22} \\
    b'' &= \frac{1 - (b')^2}{b} - \frac{a'b'}{a} + b'f' + b. \label{eq:rajan23}
\end{align}
for functions $a,b,f: [0, \infty) \to \R$ subject to initial conditions \cite[Lem.~2.1]{Rajan}
\[
    a(0) = 0,\ \ a'(0) = 1,\ \ b(0) = b_0 > 0,\ \ b'(0) = 0,\ \ f(0) = f'(0) = 0
\]
and $f''(0) = f_0 < 0$ to make use of Lem.~4.2 \cite{Rajan}. By Thm.~5.1, there exists unique solutions $a,b,f$. Let
\begin{align*}
    \alpha &\coloneq \lim_{r \to \infty} a'(r), \\
    \beta &\coloneq \lim_{r \to \infty} b'(r).
\end{align*}
Lems.~6.12--6.13 imply that the limits converge and $\alpha, \beta > 0$. We are interested in bounding $\beta$ as it determines the scalar curvature of the Ricci flow.

\section{Main Proofs}

\begin{theorem} \label{thm:main-ricci}
    Let $\lambda \coloneq \sqrt{-2f_0}$ and $c \coloneq 1 + b_0^{-2}$. If $\lambda^2 \ge 8c$, then
    \begin{equation}
        0 < \beta^2 \le \frac{b_0^2 + 1}{\lambda^2 - 4c}.
    \end{equation}
\end{theorem}

The above theorem implies that for some $C > \sqrt{b_0^2 + 1}$, we have
\[
    \beta^2 \le \frac{b_0^2 + 1}{\lambda^2 - 4c} < \frac{C^2}{\lambda^2} \implies \beta < C / \lambda,
\]
which further implies that once $f_0$ is sufficiently negative, the corresponding asymptotic cones have strictly positive scalar curvature away from their vertices.

\newpar

The following lemma is critical for proving Thm.~\ref{thm:main-ricci}.

\begin{lemma} \label{lem:monotone-q}
    Let $\lambda$ and $c$ be defined as in Thm.~\ref{thm:main-ricci}. If $\lambda^2 \ge 8c$, then \[ Q(r) \coloneq \frac{1 + [b(r)]^2}{\lambda^2 - 2f(r) - 4c} \] is strictly decreasing on $(0, \infty)$.
\end{lemma}

\begin{proof}
    Let 
    \begin{equation}
        H \coloneq \frac{a'}{a} + 2\frac{b'}{b} - f'.
    \end{equation}
    By Lems.~4.1--4.2 \cite{Rajan}, we have $a', b' > 0$ and $f', f'' < 0$ on $(0, \infty)$. Therefore, $b(r) \ge b_0$ and $f(r) \le 0$. By Eq.~\eqref{eq:rajan23}, we have
    \begin{align}
        b'' + Hb' &= \frac{1 - (b')^2}{b} - \frac{a'b'}{a} + b'f' + b + \left( \frac{a'}{a} + 2\frac{b'}{b} - f' \right) b' \notag \\ &= \frac{1 + (b')^2}{b} + b. \label{eq:gpt1}
    \end{align}
    On the other hand, Lem.~3.1 \cite{Rajan} gives
    \begin{equation} \label{eq:raj-lem31}
        \Delta f - \abs{\nabla f}^2 - 2f = 2f_0 = -\lambda^2.
    \end{equation}
    Since the Riemannian volume density of Ricci flows is proportional to $ab^2$, we have
    \begin{equation}
        \Delta f = \frac{1}{ab^2} \frac{\partial}{\partial r} (ab^2 f') = f'' + \left( \frac{a'}{a} + 2\frac{b'}{b} \right) f'.
    \end{equation}
    Substitute $\Delta f$ and $|\nabla f|^2 = (f')^2$. Then we can rewrite Eq.~\eqref{eq:raj-lem31} as follows:
    \begin{gather}
        f'' + \left( \frac{a'}{a} + 2\frac{b'}{b} \right) f' - (f')^2 - 2f = -\lambda^2 \\
        \implies f'' + Hf' = -\lambda^2 + 2f. \label{eq:gpt2}
    \end{gather}
    Since $H \ge -f' > 0$ and $f'' < 0$, by Eq.~\eqref{eq:gpt2}, we can obtain
    \begin{equation} \label{eq:aqw1}
        0 < (f')^2 = (-f')^2 \le -Hf' = f'' + \lambda^2 - 2f < \lambda^2 - 2f.
    \end{equation}
    Let
    \begin{align}
        P &\coloneq \lambda^2 - 2f - 4c, \label{eq:qp} \\
        Y &\coloneq 1 + b^2, \label{eq:qy}
    \end{align}
    where $Q = Y/P$. Our assumption $\lambda^2 \ge 8c$ gives the following relations:
    \begin{equation} \label{eq:aqw2}
        P = \lambda^2 - 2f - 4c \ge 4c - 2f \ge 4c = 4 (1 + b_0^{-2}) > 0.
    \end{equation}
    Therefore, from Eqs.~\eqref{eq:aqw1} and \eqref{eq:aqw2}, we can obtain the following
    \begin{equation} \label{eq:vbound}
        \frac{(f')^2}P < \frac{\lambda^2 - 2f}{\lambda^2 - 2f - 4c} = 1 + \frac{4c}{\lambda^2 - 2f - 4c} \le 2.
    \end{equation}
    
    \newpar

    To show that $Q$ never stops decreasing, we attempt to compute $Q''$ at the critical point of $Q$. Compute derivatives of $Y$ and $P$ given in Eqs.~\eqref{eq:qp} and \eqref{eq:qy}. We have
    \begin{equation} \label{eq:pypy}
        P' = -2f',\ \ Y' = 2bb',\ \ P'' = -2f'',\ \ Y'' = 2[bb'' + (b')^2].
    \end{equation}
    Then
    \begin{align}
        Y'' + HY' = 2(b')^2 + 2b(b'' + Hb') &= 2(b')^2 + 2b \left( \frac{1 + (b')^2}{b} + b \right) \notag \\
        &= 2b^2 + 2 + 4(b')^2 \notag \\
        &= 2Y + 4(b')^2 \label{eq:lala7}
    \end{align}
    and
    \begin{equation}
        P'' + HP' = -2(f'' + Hf') = -2(-\lambda^2 + 2f) = 2(P + 4c)
    \end{equation}
    Recall that $Y = QP$. When $Q' = 0$, we have
    \begin{gather}
        Y' = (QP)' = QP' \label{eq:lala5} \\
        Y'' = (QP' + Q'P)' = QP'' + Q''P. \label{eq:lala6}
    \end{gather}
    In Eq.~\eqref{eq:lala5}, let us replace $Y'$ and $P'$ by the expression given in Eq.~\eqref{eq:pypy}:
    \begin{equation}
        bb' = -Qf'. \label{eq:bbqf}
    \end{equation}
    From Eq.~\eqref{eq:lala6}, we can obtain
    \begin{equation} \label{eq:lala8}
        Y'' + HY' = QP'' + Q''P + HQP' = Q''P + Q(P'' + HP').
    \end{equation}
    Combining results in Eqs.~\eqref{eq:lala7} and \eqref{eq:lala8}, we have
    \begin{equation}
        PQ'' = 2Y + 4(b')^2 - Q \cdot 2(P + 4c) = 4(b')^2 - 8Qc \label{eq:b886}
    \end{equation}
    From Eq.~\eqref{eq:bbqf}, we can obtain
    \begin{equation}
        (b')^2 = \frac{Q^2(f')^2}{b^2} = \frac{QY}{P} \frac{(f')^2}{b^2} = \frac QP \cdot (f')^2 \cdot \frac{b^2 + 1}{b^2}
        = \frac{Q(f')^2 (1 + b^{-2})}{P}. \label{eq:b887}
    \end{equation}
    Finally, by Eqs.~\eqref{eq:b886}--\eqref{eq:b887}, we have
    \begin{equation}
        Q'' = \frac{4(b')^2 - 8Qc}{P} = \frac{4Q}P \left( \frac{(f')^2 (1 + b^{-2})}{P} - 2c \right).
    \end{equation}
    Eq.~\eqref{eq:vbound} gives $0 < (f')^2/P < 2$ and $b \ge b_0$ implies $1 + b^{-2} \le 1 + b_0^{-2} = c$. Therefore,
    \begin{equation}
        \frac{(f')^2}{P} (1 + b^{-2}) < 2c,
    \end{equation}
    which implies $Q'' < 0$. This implies that $Q'$ can never rise above $0$.

    \newpar

    Now, we will show that the initial value $Q'$ is negative. Consider the initial conditions and the oddness of $a$ and the evenness of $b$ and $f$, we can write
    \[
        a(r) = r + \mathcal O(r^3),\ \ b(r) = b_0 + \tfrac12 b''(0) r^2 + \mathcal O(r^4),\ \ f(r) = \tfrac12 f_0 r^2 + \mathcal O(r^4).
    \]
    Then
    \begin{align}
        H &= \frac{a'}{a} + 2\frac{b'}{b} - f' \\
        &= \frac{1 + \mathcal O(r^2)}{r + \mathcal O(r^3)} + 2 \frac{b''(0)r + \mathcal O(r^3)}{b_0 + \mathcal O(r^2)} + \mathcal O(r) \\
        &= \frac1r + \mathcal O(r).
    \end{align}
    Therefore, we have
    \begin{align}
        b'' + Hb' &= b''(0) + \left( \frac1r + \mathcal O(r) \right) (b''(0) r + \mathcal O(r^3)) \\
        &= 2b''(0) + \mathcal O(r^2). \label{eq:bbyahu1}
    \end{align}
    On the other hand, by Eq.~\eqref{eq:gpt1}, we can derive the following relations
    \begin{align}
        b'' + Hb' = \frac{1 + (b')^2}{b} + b &= \frac{1 + \mathcal O(r^2)}{b_0 + \mathcal O(r^2)} + b_0 + \mathcal O(r^2) \\
        &= b_0 + \frac1{b_0} + \mathcal O(r^2). \label{eq:bbyahu2}
    \end{align}
    Therefore, $2b''(0) = b_0 + b_0^{-1}$. It follows that
    \begin{gather*}
        Y'' = 2[bb'' + (b')^2] = 2b_0 b''(0) + \mathcal O(r^2) \\
        \implies Y''(0) = 2b_0 b''(0) = b_0 (b_0 + b_0^{-1}) = b_0^2 + 1 = Y(0)
    \end{gather*}
    and $P''(0) = -2f_0 = \lambda^2$. Also,
    \begin{gather*}
        Y' = 2bb' = 2b_0(b''(0)r + \mathcal O(r^3)) \implies Y'(0) = 0, \\
        P' = -2(f_0 r + \mathcal O(r^3)) \implies P'(0) = 0.
    \end{gather*}
    Therefore,
    \begin{equation}
        Q'(0) = \frac{Y'(0) P(0) - Y(0) P'(0)}{[P(0)]^2} = 0
    \end{equation}
    and
    \begin{gather}
        2PP'Q' + P^2 Q'' = (P^2 Q')' = Y''P - YP'' = (Y'P - YP')' \\
        \implies Q''(0) = \frac{Y''(0) P(0) - Y(0) P''(0)}{[P(0)]^2} = \frac{(b_0^2 + 1)[(\lambda^2 - 4c) - \lambda^2]}{(\lambda^2 - 4c)^2} = -\frac{(b_0^2 + 1)(4c)}{(\lambda^2 - 4c)^2} < 0.
    \end{gather}
    Hence, $Q' < 0$ for small positive $r$.
\end{proof}

With Lem.~\ref{lem:monotone-q}, we can prove Thm.~\ref{thm:main-ricci} more easily.

\begin{proof}[Proof of Thm.~\ref{thm:main-ricci}]
    Lem.~6.10 \cite{Rajan} implies
    \begin{equation}
        f'(r) = -r + \mathcal O(1) \implies f(r) = -\tfrac12 r^2 + \mathcal O(r),
    \end{equation}
    and Lem.~6.13 implies that $\lim_{r \to \infty} b(r) / r = \beta$. Recall Eqs.~\eqref{eq:qp}--\eqref{eq:qy}. We have
    \begin{gather}
        P = \lambda^2 - 4c - 2 (-\tfrac12 r^2 + \mathcal O(r)) = r^2 + \mathcal O(r) \implies P/r^2 \to 1, \\
        Y/r^2 = (b/r)^2 + 1/r^2 \to \beta^2,
    \end{gather}
    as $r \to \infty$. Therefore,
    \begin{equation}
        0 < \beta^2 = \lim_{r \to \infty} \frac{Y/r^2}{P/r^2} = \lim_{r \to \infty} \frac YP = \lim_{r \to \infty} Q(r) \le Q(0) = \frac{1 + b_0^2}{\lambda^2 - 4c},
    \end{equation}
    where $Q$ is defined as in Lem.~\ref{lem:monotone-q} the inequality is given by the lemma. We are done.
\end{proof}

\printbibliography

\end{document}